\documentclass[answers]{exam}
\usepackage{../../mypackages}
\usepackage{../../macros}

% Pour des colonnes qui prennent toute la largeur

\newcolumntype{L}{>{\raggedright\arraybackslash}p{0.48\linewidth}}

% Macro pour exercice avec pointillés automatiques
\newcommand{\ex}[1]{#1\;=\;\dotfill}

% Espace vertical plus grand entre les lignes des tableaux
\renewcommand{\arraystretch}{2.2}

\title{Calcul mental - Erreurs typiques}
\author{N. Bancel}
\date{19 novembre 2025}

% Mise en forme des solutions en bleu
\SolutionEmphasis{\color{blue}}
\renewcommand{\solutiontitle}{\noindent}

\begin{document}

\textbf{Collège Lycée Suger}
\hfill
\textbf{Mathématiques} \\

\textbf{Année 2025-2026}
\hfill
\textbf{Classe de 6ème} \par

{\let\newpage\relax\maketitle}

\begin{center}
  \textbf{\textcolor{red}{Ce document fournit des exercices de calcul mental qui correspondent aux catégories que vous avez apprises jusque là. Je vous recommande de régulièrement reprendre ces exercices à l'oral, et de vous tester en vous chronométrant et en cachant les réponses}} \\
\end{center}

\section{Fiches de calcul mental}

% ======================================================
\subsection*{Addition simple}
\begin{center}
\begin{tabular}{|L|L|}
\hline
\ex{$7+5$} & \ex{$59+18$} \\
\ex{$12+9$} & \ex{$63+29$} \\
\ex{$18+7$} & \ex{$74+36$} \\
\ex{$27+15$} & \ex{$87+45$} \\
\ex{$38+24$} & \ex{$46+27$} \\
\hline
\end{tabular}
\end{center}

% ======================================================
\subsection*{Soustraction simple}
\begin{center}
\begin{tabular}{|L|L|}
\hline
\ex{$12-5$} & \ex{$73-18$} \\
\ex{$18-9$} & \ex{$82-29$} \\
\ex{$25-7$} & \ex{$94-36$} \\
\ex{$37-15$} & \ex{$105-45$} \\
\ex{$49-24$} & \ex{$58-27$} \\
\hline
\end{tabular}
\end{center}

\subsection*{Multiplication simple}
\begin{center}
\begin{tabular}{|L|L|}
\hline
\ex{$8 \times 7$} & \ex{$6 \times 6$} \\
\ex{$11 \times 6$} & \ex{$9 \times 7$} \\
\ex{$9 \times 8$} & \ex{$6 \times 4$} \\
\ex{$4 \times 7$} & \ex{$7 \times 3$} \\
\ex{$6 \times 3$} & \ex{$9 \times 4$} \\
\hline
\end{tabular}
\end{center}


% ======================================================
\subsection*{Additions avec 99, 199, 299…}
\begin{center}
\begin{tabular}{|L|L|}
\hline
\ex{$23+99$} & \ex{$142+399$} \\
\ex{$45+99$} & \ex{$178+399$} \\
\ex{$58+199$} & \ex{$253+499$} \\
\ex{$76+199$} & \ex{$367+599$} \\
\ex{$87+299$} & \ex{$125+299$} \\
\hline
\end{tabular}
\end{center}

% ======================================================
\subsection*{Additions avec 101, 201, 1001…}
\begin{center}
\begin{tabular}{|L|L|}
\hline
\ex{$73+101$} & \ex{$675+701$} \\
\ex{$97+101$} & \ex{$19829+801$} \\
\ex{$32+201$} & \ex{$21+901$} \\
\ex{$84+301$} & \ex{$132+1001$} \\
\ex{$1530+401$} & \ex{$899+501$} \\
\hline
\end{tabular}
\end{center}

% ======================================================
\subsection*{Soustractions avec 99, 199, 299…}
\begin{center}
\begin{tabular}{|L|L|}
\hline
\ex{$123-99$} & \ex{$542-399$} \\
\ex{$145-99$} & \ex{$678-399$} \\
\ex{$258-199$} & \ex{$753-499$} \\
\ex{$276-199$} & \ex{$967-599$} \\
\ex{$387-299$} & \ex{$425-299$} \\
\hline
\end{tabular}
\end{center}

% ======================================================
\subsection*{Soustractions avec 101, 201, 1001…}
\begin{center}
\begin{tabular}{|L|L|}
\hline
\ex{$223-101$} & \ex{$842-501$} \\
\ex{$245-101$} & \ex{$978-501$} \\
\ex{$358-201$} & \ex{$1153-701$} \\
\ex{$476-201$} & \ex{$1367-1001$} \\
\ex{$587-301$} & \ex{$725-301$} \\
\hline
\end{tabular}
\end{center}

% ======================================================
\subsection*{Multiplication par 10, 100, 1000 - Simples}
\begin{center}
\begin{tabular}{|L|L|}
\hline
\ex{$7\times10$} & \ex{$59\times100$} \\
\ex{$9\times10$} & \ex{$12\times1000$} \\
\ex{$14\times10$} & \ex{$27980\times1000$} \\
\ex{$25\times10$} & \ex{$804\times1000$} \\
\ex{$38\times100$} & \ex{$461\times100$} \\
\hline
\end{tabular}
\end{center}

\subsection*{Multiplication par 10, 100, 1000 - Avec virgules}
\begin{center}
\begin{tabular}{|L|L|}
\hline
\ex{$0.7\times10$} & \ex{$5.09\times100$} \\
\ex{$0.0000237\times1000$} & \ex{$9.18\times1000$} \\
\ex{$34.18\times1000$} & \ex{$0.0192017\times100$} \\
\ex{$3.681\times100$} & \ex{$88.6\times1000$} \\
\ex{$3.1\times10000$} & \ex{$46.1\times100$} \\
\hline
\end{tabular}
\end{center}

% ======================================================
\subsection*{Division par 10, 100, 1000}
\begin{center}
\begin{tabular}{|L|L|}
\hline
\ex{$73\div10$} & \ex{$5989\div100$} \\
\ex{$99\div10$} & \ex{$12\div1000$} \\
\ex{$146\div10$} & \ex{$270\div1000$} \\
\ex{$254\div10$} & \ex{$84162\div1000$} \\
\ex{$3811\div100$} & \ex{$4671\div100$} \\
\hline
\end{tabular}
\end{center}

% ======================================================
\subsection*{Priorités de calcul}
\begin{center}
\begin{tabular}{|L|L|}
\hline
\ex{$2 + 3\times4$} & \ex{$(9-3)\times(4+2)$} \\
\ex{$5 + 2\times6$} & \ex{$2 + (3+4)\times2$} \\
\ex{$(5+3)\times2$} & \ex{$(12-4)\times(3+1)$} \\
\ex{$12 - 3\times2+(5+1)$} & \ex{$4 + 2\times(5-1)$} \\
\ex{$(12-3)\times2+4$} & \ex{$3\times(7-2)$} \\
\ex{$[(1+7)\times (3+1) - 2] \times2$} & \ex{$[(8-2)\times2 - 3\times3] - (2 + 1) \times 1$}  \\
\hline
\end{tabular}
\end{center}

% ======================================================
\subsection*{Multiplication par 4}
\begin{center}
\begin{tabular}{|L|L|}
\hline
\ex{$6\times4$} & \ex{$23\times4$} \\
\ex{$7\times4$} & \ex{$27\times4$} \\
\ex{$9\times4$} & \ex{$34\times4$} \\
\ex{$12\times4$} & \ex{$19\times4$} \\
\ex{$14\times4$} & \ex{$17\times4$} \\
\hline
\end{tabular}
\end{center}

% ======================================================
\subsection*{Division par 4}
\begin{center}
\begin{tabular}{|L|L|}
\hline
\ex{$24\div4$} & \ex{$92\div4$} \\
\ex{$28\div4$} & \ex{$104\div4$} \\
\ex{$36\div4$} & \ex{$132\div4$} \\
\ex{$48\div4$} & \ex{$52\div4$} \\
\ex{$68\div4$} & \ex{$76\div4$} \\
\hline
\end{tabular}
\end{center}

% ======================================================
\subsection*{Regroupements astucieux (multiplication)}
\begin{center}
\begin{tabular}{|L|L|}
\hline
\ex{$0.25\times78\times4$} & \ex{$0.25\times67\times8$} \\
\ex{$0.2\times45\times10$} & \ex{$0.2\times130\times100$} \\
\ex{$0.25\times17\times4$} & \ex{$0.2\times225\times5$} \\
\ex{$0.5\times84\times2$} & \ex{$0.5\times78\times4$} \\
\ex{$0.5\times246\times4$} & \ex{$0.5\times315\times2$} \\
\ex{$0.1\times0.72\times1000$} & \ex{$0.1\times100\times158$} \\
\ex{$0.25\times44\times4$} & \ex{$0.2\times175\times5$} \\
\ex{$0.5\times146\times2$} & \ex{$0.1\times390\times100$} \\
\hline
\end{tabular}
\end{center}

% ======================================================
\subsection*{Regroupements astucieux (addition)}

\begin{center}
\begin{longtable}{|L|L|}
\hline
\ex{$7 + 3 + 12$} & \ex{$37 + 3 + 40$} \\
\ex{$8 + 2 + 17$} & \ex{$48 + 2 + 50$} \\
\ex{$25 + 9 + 1$} & \ex{$63 + 17 + 3$} \\
\ex{$18 + 6 + 4$} & \ex{$24 + 6 + 30$} \\
\ex{$13 + 7 + 20$} & \ex{$45 + 15 + 5$} \\
%\hline
\ex{$14 + 6 + 11 + 9$} & \ex{$30 + 45 + 5 + 15$} \\
\ex{$8 + 22 + 4 + 6$} & \ex{$70 + 10 + 8 + 12$} \\
\ex{$12 + 8 + 14 + 6$} & \ex{$50 + 25 + 12 + 8$} \\
\ex{$33 + 7 + 3 + 17$} & \ex{$90 + 4 + 6 + 10$} \\
\ex{$5 + 15 + 25 + 8 + 2$} & \ex{$30 + 40 + 13 + 7$} \\
%\hline
\ex{$15 + 2.7 + 0.3 + 8$} & \ex{$7.8 + 12 + 0.2 + 10$} \\
\ex{$20 + 4.6 + 0.4 + 5$} & \ex{$25 + 8.9 + 0.1 + 6$} \\
\ex{$7 + 12.5 + 0.5 + 3$} & \ex{$6 + 23.75 + 0.25 + 12$} \\
\ex{$18 + 9.3 + 0.7 + 4$} & \ex{$14.6 + 25 + 0.4 + 6$} \\
\ex{$40 + 5.8 + 0.2 + 9$} & \ex{$10 + 19.9 + 0.1 + 8$} \\
\hline
\end{longtable}
\end{center}

\end{document}
